% Generated draft: requires mathematical and source review.
\documentclass[11pt]{article}
\usepackage[margin=25mm]{geometry}
\usepackage{fontspec}
\setmainfont{Latin Modern Roman}
\usepackage{amsmath,amssymb,mathtools,unicode-math}
\setmathfont{Latin Modern Math}
\usepackage{enumitem,needspace,xurl,hyperref}
\setlength{\parindent}{0pt}
\setlength{\parskip}{5pt}
\setlength{\emergencystretch}{3em}
\newcommand{\N}{\mathbb{N}}
\newcommand{\Z}{\mathbb{Z}}
\newcommand{\Q}{\mathbb{Q}}
\newcommand{\R}{\mathbb{R}}
\newcommand{\C}{\mathbb{C}}
\newcommand{\F}{\mathbb{F}}
\newcommand{\A}{\mathbb{A}}
\newcommand{\PP}{\mathbb{P}}
\newcommand{\E}{\mathbb{E}}
\newcommand{\eps}{\varepsilon}
\newcommand{\dd}{\,\mathrm{d}}
\begin{document}
\section*{UnsolvedMath ID: 30004232}
\textbf{Source problem code:} catalogue entry
\subsection{Definitions and mathematical statement}
Let \(n\) be a positive integer and write \([n]=\{1,\ldots,n\}\); the source does not explicitly state the range of \(n\). Let \(G\) be a bipartite graph with parts \(X=[n]\) and \(Y\). For every nonempty \(S\subseteq[n]\), define
\[
X_S=\bigl|\{y\in Y:N_G(y)=S\}\bigr|,
\]
where \(N_G(y)\) is the neighborhood of \(y\). These counts must be finite for their numerical use below; the source does not explicitly impose this condition. Vertices with empty neighborhood do not enter the displayed definitions.

For the polynomial formulation, regard the symbols \(X_S\) as independent indeterminates and work in
\[
R_n=\Q[X_S:\varnothing\ne S\subseteq[n]].
\]
This rational coefficient ring accommodates the displayed coefficients; the source does not specify the coefficient ring for its requested basis. Define \(M\in\operatorname{Mat}_{n\times n}(R_n)\) by the formulas exactly as supplied:
\[
M_{ii}=\sum_{\substack{\varnothing\ne S\subseteq[n]\\i\in S}}|S|,
\qquad
M_{ij}=\sum_{\substack{\varnothing\ne S\subseteq[n]\\i\in S,\ j\notin S}}\frac{X_S}{|S|}
\quad(i\ne j).
\]
In particular, the supplied diagonal formula contains no factor \(X_S\). Whether this omission is intentional is unresolved; no such factor has been inserted here. Let
\[
D=\operatorname{diag}(M_{11},\ldots,M_{nn}),\qquad P_n=\det M,\qquad Q_n=\det D.
\]
The relabeling action is the action of \(\mathfrak S_n\) on \(R_n\) defined by \(\sigma\cdot X_S=X_{\sigma(S)}\). The source identifies \(P_n\) and \(Q_n\) as invariant under this action. This action permutes subset labels through permutations of \([n]\); it does not assert symmetry under arbitrary permutations of all the indeterminates.

\textbf{Requested mathematical target.} Find a natural basis from symmetric-function theory in which to express \(P_n\) and \(Q_n\). The source does not specify the ambient space having this basis, the meaning of naturality, or whether a compatible construction for all \(n\) is required. Thus this is a basis-construction question whose intended precise formulation remains to be clarified.

The original statement also asks whether a proposed sum-of-squares refinement of the difference of these determinants is valid. No formula for that refinement, order of subtraction, domain of validity, or allowed square terms is supplied. The background describes a stronger nonnegative-square decomposition, but supplies no additional defining expression. This additional target cannot be stated as a precise conjecture from the supplied record.
\subsection{Short English statement}
The problem asks for a natural symmetric-function basis for two determinant polynomials invariant under relabeling one part of a bipartite graph. The diagonal formula is retained exactly as supplied, although its lack of an \(X_S\) factor requires review. An additional sum-of-squares question is mentioned without its proposed formula. The supplied dated curation labels the remaining questions open and reports the determinant inequality proved; these assessments are not independently verified here.
\subsection{Sources}
\begin{itemize}
\item Local problem registry \texttt{data/unsolvedmath/problems.json}, permanent UnsolvedMath ID \(30004232\), exact source problem code \texttt{catalogue entry}, title ``Symmetric Function Bases for Bipartite Determinant Inequalities.'' This draft uses the supplied complete source record. The recorded proposed year is 2019; no individual proposer is supplied.
\item Supplied report citation: \emph{Geometric, Algebraic, and Topological Combinatorics} (2020), \emph{Oberwolfach Reports}, pp.~2395--2472, DOI 10.4171/owr/2019/39. \url{https://doi.org/10.4171/owr/2019/39}. The record attributes the problem to workshop open-problem discussions collected in this report.
\item Supplied literature source: ``Ho, determinant inequality.'' \url{https://arxiv.org/abs/2603.17997}. No fuller author identification or bibliographic details are supplied.
\item Literature assessment checked at \texttt{2026-08-22}: the supplied curation labels the problem open and reports that the determinant inequality is proved, while the requested symmetric-function basis and stronger nonnegative-square decomposition are not established by the available literature. This is supplied dated curation, not an independently verified assessment; the record requests independent verification of time-sensitive status and cited claims.
\item The registry assigns statement status \texttt{corrected\_verified} and reports that the official Oberwolfach report was checked directly and that definitions were restored in the clean statement. These are supplied verification claims, not verification performed for this draft.
\end{itemize}
\end{document}
